\section{本讲小结}
功是力与作用点位移的积分，冲量是力的时间积分；前者决定动能变化，后者决定动量变化。计算功率时同样应使用受力作用点的速度。弹簧、滑轮、浮力和连续喷流中的几何或流量关系，必须先与所选积分变量对应。

能量式给出速度以后，还应检查接触力、绳张力、路径上的首个零动能点及事件顺序。多阶段问题在分离、绳松弛和碰撞时切换方程；用一个始末等式代替全程可达性分析，容易得到形式成立却无法实现的状态。
